\section{Задание}

\section{Задание}

\begin{enumerate}[label=\textbf{Task \arabic*:}, leftmargin=*]

\item \textit{[2 балла]}

Выберите датасет и задачу для него (классификация) --- с Kaggle или \href{https://archive.ics.uci.edu/datasets}{UCI Machine Learning Repository}.

Рекомендации:
\begin{itemize}
    \item количество признаков --- не менее 10;
    \item количество объектов --- не менее 100;
    \item возраст датасета --- не старее 2010 года;
    \item пользуйтесь датасетом из предыдущих лабораторных работ, если он подходит под условия выше.
\end{itemize}

Затем выполните следующие действия:
\begin{itemize}
    \item Выберите (или сгенерируйте сами) признак, который будет целевым для бинарной классификации.
    \item Используйте LogisticRegression из sklearn. Задайте несколько параметров метода самостоятельно, поясните их назначение и причину выбора конкретных значений.
    \item Используйте несколько метрик, чтоб оценить результаты бинарной классификации (метрики, доступные в sklearn).
\end{itemize}

\item \textit{[2 балла]}

Выполните бинарную классификацию аналогично предыдущему пункту с использованием SVM из sklearn. Используйте \texttt{kernel=‘linear’}, сами подберите параметр регуляризации С и ответьте на вопрос:
\begin{itemize}
    \item Может ли SVM учесть несбалансированность классов? Если да, то как? Если нет, то почему?
\end{itemize}

\item \textit{[6 баллов]}

\begin{enumerate}
    \item Многоклассовая классификация.[\textcolor{DarkOrange}{2 балла}]
    \begin{itemize}
        \item Дайте определения для подходов: микро, макро, взвешенный, пообъктный.
        \item Выполните расчёты вручную для слайда 43 (на примере любых 3 промежуточных AUC ROC, поясните, как из промежуточных характеристик получаются 4 финальных).
        \item В рамках детализации \textbf{оцифруйте} ROC: в Excel, отрисуйте линиями в Word или отрисуйте отдельными отрезками Python-кодом (не используйте готовые функции для ROC).
    \end{itemize}
    \item Для одного из кейсов (одного класса):[\textcolor{DarkOrange}{2 балла}]
    \begin{itemize}
        \item Введите три порога бинаризации так, чтоб получились различные Accuracy, Recall, Precision.
        \item Вычислите эти метрики вручную и сопроводите визуализацией расчёта для Recall и Precision (см. пример из лекции с окружностями разных цветов и радиусов).
        \item По трём парам (Precision, Recall) постройте эскиз PR-кривой.
    \end{itemize}
    \item Сделайте выводы по результатам расчётом AUC ROC для многоклассовой классификации:[\textcolor{DarkOrange}{1 балл}]
    \begin{itemize}
        \item какой вариант лучше;
        \item почему;
        \item объясните падение AUC ROC при переходе к многоклассовым.
    \end{itemize}
    \item В каких кейсах/доменах следует применять те или иные варианты? Запишите выводы. [\textcolor{DarkOrange}{1 балл}]
\end{enumerate}

\item \textit{[2 балла]}

Исследуйте вопрос выбросов/новизны для данных в вашем датасете с использованием OneClassSVM. Дайте краткую теоретическую характеристику, используйте готовые решения для получения результатов, сделайте выводы.

Материалы для изучения указаны в списке литературы.\cite{sklearn_outlier_detection,dyakonov_anomaly_search}

\end{enumerate}